%% ====================================================================
%%  Chapter 4 — Implementation
%% ====================================================================
\chapter{Implementation}
\label{ch:implementation}

This chapter describes how the proposed methods were implemented. It presents the codebase, together with
the tools used to precompute the teacher outputs, train the student models, evaluate their robustness,
test their generalization, and make the experiments repeatable across different random seeds.

\section{Codebase and Training Infrastructure}
\label{sec:impl:infra}

All of the models were implemented in PyTorch. The codebase is organized as a core package that contains
the DLoc model generators, a data loader, and evaluation and training utilities; a models package that
contains the student model classes (MobileNetV2, TinyCNN, MobileNetV2-UNet, and the Mamba variant); and
an evaluation package that contains the scripts used for benchmarking the models and generating the
figures in this manuscript.

The models were trained on rented NVIDIA H100 GPUs (Vast.ai). Lighter benchmarking and figure generation
were carried out on a local NVIDIA RTX 3060 GPU. All of the trained models, training logs, evaluation
files, and precomputed teacher outputs are archived on the Hugging Face Hub---the model artifacts under
\texttt{JustAGeek/dloc-code} and the WILD measurement files under the dataset repository
\texttt{JustAGeek/dloc-wild-fig10b}---so that every reported number can be regenerated from the released
weights and data.

The code was written to allow for repeatable experiments, and one practical lesson is worth recording.
The localization error of Equation~\eqref{eq:method:error} is computed as a per-sample loop over the
predicted and ground-truth locations. In its original implementation, NumPy's threaded backend was
invoked for each sample; on a 192-core machine this caused the math libraries to spin up their entire
thread pool for each tiny operation, pinning many cores while doing little useful work. Constraining the
math libraries to a single thread (via the \texttt{OMP\_NUM\_THREADS} family of environment variables)
removes this overhead entirely and is required for the evaluation scripts to run at full speed.

\section{Teacher Output Precomputation}
\label{sec:impl:teacher}

To enable training of the students in distilled mode, the teacher output $y_T$ in
Equation~\eqref{eq:method:distill} is precomputed for every example in the training and test sets. These
outputs are cached to disk (approximately 4.09\,GB for the training split and 0.99\,GB for the test
split). During distilled training, the students are trained to minimize the distance between their
outputs and these cached teacher outputs, so the teacher model is never invoked during student training.
This precomputation of the teacher's output is the only change required to make distilled training
tractable across many students and seeds.

\section{Student Training}
\label{sec:impl:training}

The convolutional students (MobileNetV2 and TinyCNN) were trained for 50 epochs with the Adam optimizer
at a learning rate of $10^{-4}$ and a batch size of 64. The MobileNetV2-UNet model was trained for 120
epochs, also with Adam, under one of two different learning-rate schedules. For the cosine-annealing
schedule, the learning rate was annealed from its initial value of $10^{-4}$ to zero over the course of
the 120 epochs. For the step-decay schedule, the learning rate was halved every 20 epochs. The DLoc
baseline is trained using Adam ($\beta_1 = 0.5$) with a StepLR schedule ($\gamma = 0.9$) whose step sizes are 50, 20, and 50
epochs for the encoder, location decoder, and offset decoder, respectively. Each model was trained in
both its standalone and distilled modes; in distilled mode $\alpha$ was set to 0.7 and $\beta$ to 0.3.
For each student, the checkpoint with the minimum test-set error was selected as the best checkpoint, and
per-epoch metrics for both the training and test sets were recorded for every training run.

\section{Robustness Test Harness}
\label{sec:impl:robustness}

Robustness is evaluated by applying signal degradations to the CSI measurements at test time, without
retraining the models. For each model, the signal is corrupted in six different ways: by dropping the
signal from one of the four APs (either individually or at random), by adding Gaussian noise of
$\sigma \in \{0.1, 0.2\}$, by globally attenuating the signal by 50\%, by attenuating the signal from one
AP to 30\% of its original strength, and by applying Gaussian blur to the CSI signal with a $5 \times 5$
kernel and $\sigma = 1$. Each model is evaluated under every condition, producing the per-condition
tables of Section~\ref{sec:res:robustness}.

\section{Cross-Session Leave-One-Session-Out Pipeline}
\label{sec:impl:crosssession}

The cross-session experiment was performed with the same protocol as the original DLoc
leave-one-session-out method. Three folds are used, each of which holds out one of the three August
sessions from the training set and trains on the remaining setups. Within each fold, the MobileNetV2-UNet
is trained according to its own recipe (Adam at $10^{-4}$ with cosine annealing over 120 epochs) for
three different random seeds. Scripts automate this process, training the models for the three seeds
within each fold and aggregating the results from each run into the mean-and-standard-deviation table of
Section~\ref{sec:res:crosssession}. Unlike the distillation experiments, the DLoc teacher model is not
retrained on these folds; the results of this experiment are directly comparable to the published results
of DLoc. Because distillation is performed at training time, it would be necessary to retrain a teacher
model for each fold, so the models were trained in standalone mode only within these cross-session
experiments. The purpose of these experiments is therefore to report the performance of the standalone
MobileNetV2-UNet model.

\section{Multi-Seed and Benchmarking}
\label{sec:impl:benchmark}

In addition to testing the models for in-distribution accuracy, cross-session accuracy, and robustness,
the DLoc baseline was trained with three different random seeds to characterize its seed-to-seed
variance. A separate script measures the computational efficiency of the models: it evaluates each
model's number of parameters, the number of floating-point operations (FLOPs) it performs on dummy
inputs, and the latency of its predictions on both CPU and GPU. These results are presented in
Section~\ref{sec:res:efficiency}. Finally, a figure-generation script regenerates the per-sample error
arrays of each model from the released checkpoints and the measurement data, and uses them to render the
cumulative-distribution, compression, cross-session, and qualitative-heatmap figures presented in
Chapter~\ref{ch:results}, so that the figures are reproducible rather than hand-made.
